% commentaire.tex — f166r, vv. 295–336 (conversion initiale des travaux antérieurs)
\comm{295–311}{}{La coryphée-prêtresse commence la prière, en prose selon l'usage : (scholie) :\citComm{πεζῇ εὐφημίᾳ χρῶντα οἱ κωμικοὶ, ἐπειδὰν εὐχὴν ἢ ψήφισμα εἰσάγωσιν.}Cf. en \italique{Oiseaux} (vers 864-888) pour une prière analogue.}
\comm{307–8}{τὸν δῆμον τῶν Ἀθηναίων καὶ τὸν τῶν γυναικῶν.}{Plaisanterie \italique{contre l'attente} (παρὰ προσδοκίαν) : « Le peuple des Athéniens et celui... des femmes », là où l'on attendait \grec{καὶ τὸν τῶν ξυμμάχων}, « et celui de ses alliés ».}
\comm{312–30}{Chant du chœur.}{Pour l'étude métrique de ce chant du chœur, l'on pourra consulter Lætitia Parker et son ouvrage de référence, \italique{The Songs of Aristophanes}\nc{Lætitia P. E. Parker, \italique{The songs of Aristophanes}, Clarendon Press Oxford (1997), pp. 406-409.}. Elle parvient à faire apparaître comme une sorte de refrain scandé $- \cup \cup - \cup \cup - \cup \cup -$\nc{Lætitia désigne ce vers comme un \italique{énoplien} (dactylo-épitrite prolongé) et Prato comme un \italique{alcméen}.} aux vers 315, 320 et 325.\\
Quoique la découverte soit extrêmement intéressante et encore à creuser, j'ai préféré pour ma part reprendre la colométrie de Prato, plus proche du manuscrit de Ravenne : «~\italien{Il Coro accoglie l'invito dell'Aralda e rivolge la sua preghiera alle divinità peculiari dell' ecclesia (maschile) escluse nella precedente invocazione in prosa e qui menzionate, non senza adattamenti, col consueto criterio selettivo, quasi per conferire un crisma di legalità maggiore alle decisioni finali.}~» (Prato.)}
\comm{330}{εὐγενεῖς γυναῖκες.}{Confirme, s'il était besoin, que seules les citoyennes athéniennes assistaient aux Thesmophories (cf. le vers 294 lorsque Mnésiloque renvoie Thratta).}
\comm{336–7}{ἢ 'πικηρυκεύεται / Εὐριπίδῃ Μήδοις ἐπὶ βλάβῃ τινὶ.}{Dans sa parodie d'ecclésie, Aristophane ne manque pas de reprendre les imprécations contre les Mèdes, qui étaient de règle comme on le sait par Isocrate (\italique{Panégyrique}, 157) :
\citComm{Καὶ πολλῶν μὲν οἱ πατέρες ἡμῶν μηδισμοῦ θάνατον κατέγνωσαν, ἐν δὲ τοῖς συλλόγοις ἔτι καὶ νῦν ἀρὰς ποιοῦνται, πρὶν ἄλλο τι χρηματίζειν, εἴ τις ἐπικηρυκεύεται Πέρσαις τῶν πολιτῶν.\\
\francais{Nos pères ont condamné à mort plusieurs citoyens pour leur adhésion au parti des Perses (μηδισμός). Encore aujourd'hui, dans nos assemblées, avant de traiter aucune affaire, on prononce des imprécations coutre celui des citoyens qui fait des ouvertures aux Mèdes.}}
L'ennemi héréditaire, pour la féminine assemblée, est remplacé -- on s'y attendait ! -- par Euripide.
Remplacé ? pas tout-à-fait ; car au côté d'Εὐριπίδῃ continuent de subsister étrangement les Μήδοις... détail qui fleure bon son interpolation : le spectateur athénien n'avait pas besoin qu'on lui souligne ainsi le trait comique...
Or nous possédons un papyrus d'Oxyrhynque\nc{\italique{The Oxyrhynchus Papyri}, IX (1912), ed. Arthur S. Dunt. Il s'agit du papyrus n° 1176, fragment 39, colonne XII 8-16. Ce papyrus date du \nieme{II} siècle ap. J.-C. et comporte un fragment d'une \italique{Vie} d'Euripide par Satyros de Callatis. Je me réfère aussi à l'édition ainsi qu'au commentaire de ce fragment par Stefan Schorn, \italique{Satyros aus Kallatis} (2002).} du \nieme{II} siècle après Jésus-Christ\nc{Quand notre manuscrit de Ravenne est beaucoup plus tardif : \nieme{IX} ou \nieme{X} s. ap. J.-C.} qui comporte une citation des vers 335-337 des \italique{Thesmophories}. Malheureusement, le vers qui nous intéresse est lacunaire :\citComm{ευριπιδιω\douteux{ι}[. .]\douteux{υ}\nc{Schorn indique que la présence d'un \douteux{c} est moins probable, et que ce n'est assurément pas un \douteux{ε}.}τ|επιβλαβη\douteux{ι}τι|νι}Hunt avait déjà remarqué en son temps que la lacune était bien insuffisante pour qu'ait pu s'y trouver Μήδοις τ᾽\nc{Ce qui n'a pas empêché les éditeurs depuis de prétendre dans leur apparat critique que ce papyrus confirmait l'addition du τ᾽ de Biset : Εὐριπίδῃ Μήδοις <τ᾽>... sauf que le τ identifié par Hunt ne pouvait en aucun cas se rattacher aux Μήδοις, absents du papyrus.}.}
